\chapter{Resultados}\label{ch:results}

Este capítulo apresenta as análises calculadas e verificadas: as superfícies de aptidão
do presente (Seção~\ref{sec:atlas-present}), o zoneamento agroambiental
(Seção~\ref{sec:zoning-results}), a mudança projetada pelo CMIP6 (Seção~\ref{sec:shift}) e a fase de
comparação com o uso atual da terra (Seção~\ref{sec:realized}). Todos os valores quantitativos provêm
das saídas verificadas do modelo. As tabelas de referência e suplementares e as figuras de
diagnóstico estão nos Anexos.

\section{Atlas de aptidão multissegmento do presente}\label{sec:atlas-present}

O mecanismo baseado em conhecimento (Seção~\ref{sec:suitability}) produziu, para cada um dos sete
segmentos, uma superfície de aptidão do presente sobre GO~e~DF a 250~m
(Figura~\ref{fig:atlas-present}). Cada superfície traz uma classe da FAO (S1/S2/S3/N) e uma faixa de
robustez de sensibilidade dos pesos de $\pm 20\%$. Três verificações confirmam o comportamento esperado
do mecanismo. Primeiro, todas as faixas de valores das bandas estão dentro dos intervalos admissíveis.
Segundo, as máscaras de cobertura da terra excluem, como esperado, as áreas urbanas de Goiânia e
Brasília e os grandes reservatórios. Terceiro, os padrões espaciais correspondem à expectativa
agronômica: a aptidão das culturas se concentra nos planaltos planos e de solo profundo, e a aptidão
para conservação cresce no terreno íngreme, nativo e ripário. O resumo quantitativo mais compacto é a
aptidão média de cada segmento por zona (Tabela~\ref{tab:zone-suit}, Seção~\ref{sec:zoning-results}),
acompanhada da parcela da terra avaliada de cada segmento que é apta (classe S2+).

\begin{figure}[htbp]
\centering
\includegraphics[width=0.78\textwidth]{fig_4_1.png}
\caption{Superfícies de aptidão do presente (0--1) para os sete segmentos. Os mapas de
classes da FAO para segmentos representativos estão na Figura~\ref{fig:fao-classes} do
Anexo.}
\label{fig:atlas-present}
\end{figure}

Soja e cana-de-açúcar pontuam mais alto nos planaltos mais frios, altos e argilosos (parcelas
S2+ de 55,8\% e 42,1\%) e caem acentuadamente em terreno íngreme ou arenoso. As outras culturas
anuais acompanham a regularidade da chuva e o calor acumulado, além da declividade e da
fertilidade (S2+ de 45,1\%). A pecuária é graduada pela produtividade de forragem, de modo que
pastagens de fato ruins caem para S3 ou N. Sua parcela S2+ é de 76,0\%, mas sua aptidão média
territorial fica \emph{abaixo} da aptidão das culturas na terra boa.

A piscicultura tem 22,0\% da sua terra avaliada na classe S2+. Essa superfície só se afasta do piso
porque o seu fator de proximidade está ancorado na água superficial \textbf{semipermanente}, e
não na permanente (Anexo~\ref{anx:derivations}).

A superfície de valor de conservação é de fato informativa (S2+ de 38,5\%). Ela agrega, de forma
compensatória, a conectividade a áreas protegidas, a heterogeneidade do terreno, o serviço de
armazenamento de carbono, o serviço de água e a marginalidade da terra. É bem graduada e identifica
uma zona cujo melhor uso absoluto é a própria conservação: a escarpa íngreme, onde a
aptidão das culturas colapsa (Seção~\ref{sec:zoning-results}, Zona~2). A geração fotovoltaica
atinge 26,9\% em S2+ e distingue as planícies mais baixas, quentes e de alta irradiância dos
planaltos mais frios. Essas planícies são favorecidas pela irradiância e pelo índice de claridade,
mas penalizadas pela redução de desempenho por calor no trimestre mais quente.

O território de GO~e~DF é climaticamente quase uniforme. Por isso, as pertinências climáticas saturam,
e a discriminação real entre terra boa e ruim vem do \textbf{solo e do terreno}. Uma decomposição da
variância espacial da aptidão por fator torna essa leitura quantitativa
(Seção~\ref{sec:factor-variance-annex} do Anexo). Quatro das sete superfícies (soja, cana-de-açúcar,
outras culturas anuais e piscicultura) dependem quase só de solo, terreno e hidrografia, com
contribuição climática próxima de zero; o fator dominante, porém, varia (na soja, a declividade
explica sozinha 93\% da variância). As outras três superfícies são exceções. A pecuária e a geração
fotovoltaica são governadas por fatores de origem climática (umidade do solo e irradiância), e a
conclusão anterior sobre o clima não vale para elas. A conservação, o segmento com mais fatores,
concentra dois terços da variância em um só deles: a conectividade a áreas protegidas, o único fator
dominante que é uma camada de governança, e não uma variável de superfície.

Essa concentração não decorre de um peso AHP enviesado a favor dos fatores dominantes. Pelo
contrário: a derivação dos pesos (Seção~\ref{sec:ahp}) desloca peso deliberadamente \emph{para} os
fatores de maior poder discriminante em cada segmento, o oposto de um viés que inflaria
artificialmente a concentração de variância. Nos três segmentos de cultivo anual, portanto, os fatores
climáticos delimitam o intervalo possível de aptidão, mas não desenham a estrutura espacial fina do
atlas, que depende do solo e do terreno. Essa dependência é retomada como limitação na
Seção~\ref{sec:limitations}.

A decomposição é calculada sobre os pixels amostrados: mostra \emph{quanto} cada fator contribui em
média, mas não \emph{onde} nem em que escala. Para isso, calculou-se o desvio-padrão local (janela
móvel) de cada banda padronizada do conjunto de atributos, agregado por tema (clima e terreno/solo) em
três raios de janela. Um deles, de 4,75~km, corresponde ao tamanho da célula nativa do TerraClimate
(Figura~\ref{fig:scale-mismatch} e Tabela~\ref{tab:scale-mismatch} do Anexo). É a escala mais
favorável ao clima, por ser a sua própria resolução nativa. Mesmo nela, a variabilidade espacial local
de terreno/solo é 41 vezes maior que a do clima ($\log_2 = 5{,}37$); nos raios mais finos (750~m e
250~m), a razão sobe para 197 e 318 vezes.

Essa uniformidade é \emph{espacial} (entre lugares, em um mesmo instante), e não \emph{temporal}
(em um mesmo lugar, ao longo do tempo). Uma superfície climática pode ser quase plana no espaço e,
ainda assim, se deslocar substancialmente no tempo, em todo o território de uma vez. É
precisamente esse deslocamento temporal que a projeção CMIP6 explora (Seção~\ref{sec:shift}).

A Figura~\ref{fig:climate-shift} torna esse deslocamento concreto para as duas entradas
climáticas mais diretamente alteradas pelo conjunto de GCM: sob SSP5-8.5, em 2051--2070, a precipitação
anual cai e a temperatura média do trimestre mais quente sobe em todo o território. O padrão espacial
de cada mapa é suave, sem a estrutura fina de bacias e microrregiões que aparece nos mapas de solo e
terreno, e permanece igualmente suave no futuro. Isso ocorre porque a mesma transformação por fatores
de mudança (Seção~\ref{sec:cmip6}) é aplicada ponto a ponto: ela desloca o território inteiro de uma
vez, sem introduzir estrutura espacial nova.

\begin{figure}[htbp]
\centering
\includegraphics[width=0.95\textwidth]{fig_4_14.png}
\caption{Entradas climáticas no presente (linha de base 1991--2020) e no futuro (SSP5-8.5,
média do conjunto em 2051--2070): precipitação anual (acima) e temperatura média do trimestre mais
quente (abaixo).}
\label{fig:climate-shift}
\end{figure}

\section{Zoneamento agroambiental}\label{sec:zoning-results}

O número de zonas, $k=4$, foi adotado por maximizar o índice de Rand ajustado medido
\emph{fora da amostra}, na metade de validação da partição espacial (Seção~\ref{sec:zoning}); o painel
completo de validade e estabilidade está na Figura~\ref{fig:cluster-validity} do Anexo. As quatro
zonas (Figura~\ref{fig:zones}) dividem o território em parcelas de 46,4\%, 10,4\%, 15,2\% e 28,1\% da
área de interesse. As médias biofísicas discriminantes por zona estão na Tabela~\ref{tab:zone-means}
do Anexo.

\begin{figure}[htbp]
\centering
\includegraphics[width=\textwidth]{fig_4_3.png}
\caption{Zonas agroambientais ($k = 4$) sobre GO~e~DF.}
\label{fig:zones}
\end{figure}

\begin{table}[htbp]
\centering
\caption{Aptidão média por segmento e por zona, com o rótulo de melhor segmento
comparativo.}
\label{tab:zone-suit}
\small
% GENERATED by tools/gen_thesis_tables.py -- do not edit by hand.
% Edit the source CSV and regenerate; see CLAUDE.md §9.
\begin{tabular}{@{}lrrrr@{}}
\toprule
Segmento & Z1 & Z2 & Z3 & Z4 \\
\midrule
Soja & 0,58 & 0,03 & 0,59 & 0,59 \\
Cana-de-açúcar & 0,52 & 0,06 & 0,47 & 0,41 \\
Outras culturas anuais & 0,54 & 0,12 & 0,46 & 0,44 \\
Piscicultura & 0,22 & 0,04 & 0,53 & 0,30 \\
Pecuária & 0,55 & 0,41 & 0,54 & 0,70 \\
Conservação & 0,42 & 0,56 & 0,45 & 0,45 \\
Geração fotovoltaica & 0,32 & 0,11 & 0,33 & 0,52 \\
\midrule
\textbf{Vocação primária} & outras anuais & conservação & piscicultura & fotovoltaica \\
\textbf{Vocação secundária} & cana & pecuária & soja & pecuária \\
\textbf{Margem} ($z$) & 0,04 & 1,31 & 0,60 & 0,05 \\
\bottomrule
\end{tabular}

\end{table}

As quatro zonas admitem uma interpretação fisiográfica direta (planalto de cultivo, escarpa
de conservação, baixada drenada e baixada quente de pastagem) e são rotuladas pelo melhor
segmento comparativo (Seção~\ref{sec:zoning}): o \emph{argmax} da aptidão normalizada em escore $z$, e
não da aptidão absoluta.

Esse rótulo, porém, só é informativo quando a distância até o segundo colocado o sustenta, e em
duas das quatro zonas não o sustenta. Por isso, a Tabela~\ref{tab:zone-suit} apresenta a vocação
primária, a secundária e a margem em escore $z$ entre as duas. Em duas zonas, a escarpa de
conservação (1,31) e a baixada drenada, voltada à piscicultura (0,60), a margem é ampla, e nelas o
rótulo primário pode ser lido como a vocação da zona. Nas outras duas, a margem é de 0,04 e 0,05, uma
ordem de grandeza menor: no planalto, as outras culturas anuais empatam com a cana-de-açúcar; na
baixada quente, a geração fotovoltaica empata com a pecuária. Essas duas zonas somam
\textbf{74,4\% do território} e devem ser lidas como pares de usos, e não como vocações únicas.

Duas consequências de $k$ ser menor que o número de segmentos merecem registro. A primeira é
que três segmentos (soja, cana-de-açúcar e pecuária) não recebem rótulo primário em nenhuma zona e
aparecem apenas como vocação secundária; a partição em quatro zonas não cobre os sete segmentos, nem
foi concebida para isso. A segunda é que a soja, apesar de ter a maior aptidão absoluta em duas das
quatro zonas, nunca é o melhor uso \emph{comparativo} de nenhuma. No planalto, porém, ela perde o
rótulo por apenas 0,04.

\begin{itemize}
  \item \textbf{Zona 1: planalto de cultivo alto e argiloso (46,4\%), vocação outras
    culturas anuais e cana-de-açúcar (margem 0,04).} A maior das quatro zonas é o núcleo
    agrícola do território: a mais alta (784~m), a mais argilosa (37,0\%) e a mais fria
    (24,5~$^\circ$C no trimestre mais quente), com o menor déficit hídrico climático anual
    (411~mm) e o maior índice de aridez, isto é, a menos seca (1,15). Concentra a maior parcela
    de soja cultivada entre as zonas (23,1\%, contra 7,9\% e 5,4\% nas duas baixadas). O rótulo
    comparativo é das outras culturas anuais, mas por apenas 0,04 em aptidão normalizada à frente da
    cana-de-açúcar. Não há, nessa zona, um segmento comparativamente dominante, e a leitura correta é
    que o planalto é bom para o cultivo anual em geral. A soja tem aqui aptidão absoluta de 0,58,
    valor alto, mas não o maior do território: as duas baixadas chegam a 0,59, com vantagem de
    milésimos para a baixada quente (0,591, contra 0,588 da drenada).
  \item \textbf{Zona 2: escarpa íngreme de conservação (10,4\%), vocação conservação
    (margem 1,31).} A menor das quatro zonas é a única com vocação comparativa inequívoca (a
    margem é mais de trinta vezes a do planalto) e também a única em que a conservação é o
    melhor uso \textbf{absoluto} (0,56), e não só comparativo. É de longe a mais íngreme
    (16,9$^\circ$, quase \emph{quatro vezes} a segunda mais íngreme, com 4,5$^\circ$), a de menor
    índice de umidade topográfica (8,8) e a mais rica em carbono orgânico (16,0~g/kg) e em água
    disponível no solo (31,9\%). A aptidão das culturas colapsa: soja, 0,03; piscicultura,
    0,04; cana-de-açúcar, 0,06, uma ordem de grandeza abaixo das demais zonas. A cobertura nativa
    chega a 69,7\%, de longe a maior do território, e a soja cultivada, a 0,1\%: na prática, o
    relevo já exclui o cultivo.
  \item \textbf{Zona 3: baixada drenada (15,2\%), vocação piscicultura e soja (margem
    0,60).} É definida quase inteiramente pela hidrografia: densidade de drenagem de 0,107,
    \textbf{quase nove vezes} a de qualquer outra zona, o maior índice de umidade topográfica
    (11,4) e a menor distância à água (8.510~m, contra 15.620--19.624~m nas demais). Daí a maior
    aptidão para piscicultura do território (0,53, contra 0,30 na segunda colocada) e a
    vantagem comparativa mais pronunciada entre as três zonas de cultivo. O uso atual, porém, é
    dominado por pastagem (58,1\%), com soja residual (7,9\%). É a divergência entre potencial e
    uso atual mais acentuada do zoneamento (Seção~\ref{sec:realized}).
  \item \textbf{Zona 4: baixada quente e arenosa de pastagem (28,1\%), vocação geração
    fotovoltaica e pecuária (margem 0,05).} É a mais quente (26,6~$^\circ$C), a mais arenosa
    (26,4\% de argila), a mais pobre em carbono orgânico (11,1~g/kg), a mais seca (aridez de 1,05)
    e a de maior déficit hídrico anual (475~mm). O calor e o solo arenoso deprimem mais as culturas
    (a cana-de-açúcar cai a 0,41, a menor das três zonas cultiváveis) do que a pastagem, que
    alcança aqui sua maior aptidão territorial (0,70). O uso atual acompanha esse padrão: 66,5\% de
    pastagem, a maior parcela entre as zonas. A margem comparativa, porém, é estreita: a geração
    fotovoltaica supera a pecuária por apenas 0,05 em aptidão normalizada, e esta zona é mais bem lida
    como de vocação dupla, fotovoltaica e pastagem. As duas leituras convergem na prática: ambos os
    usos toleram o calor e o solo arenoso que penalizam o cultivo, e é essa tolerância, e não a
    escolha entre eles, que a zona identifica.
\end{itemize}

\section{Mudança projetada pelo clima (QP3)}\label{sec:shift}

A projeção por fatores de mudança (Seção~\ref{sec:cmip6}) foi concluída para os dois cenários
(SSP2-4.5 e SSP5-8.5) e as duas janelas (2031--2050 e 2051--2070). Para cada combinação, a projeção
produz: a aptidão futura por segmento; a mudança $\Delta S$ em relação à linha de base rederivada; e
uma superfície de concordância entre modelos. Essa superfície sustenta a leitura de robustez usada ao
longo deste capítulo: ela mede, célula a célula, a fração dos cinco GCM que concorda com o
\emph{sinal} da mudança média. No caso de maior mudança projetada (outras culturas anuais, SSP5-8.5,
2051--2070), a média territorial é \textbf{0,999}. Duas
variáveis climáticas (radiação solar e umidade do solo) são mantidas na linha de base, o que suprime
parte dos sinais da geração fotovoltaica e da pecuária.

O fato de apenas o clima variar não contradiz a quase uniformidade espacial documentada na
Seção~\ref{sec:atlas-present} (Figura~\ref{fig:scale-mismatch}). Hoje, a maioria das pertinências
climáticas está \textbf{saturada} (por exemplo, a precipitação anual da soja está bem dentro da faixa
ótima em quase todo o território); por isso, o clima quase não discrimina o espaço no presente. Sob os
fatores de mudança do CMIP6, porém, a mesma variável pode sair do platô e entrar no trecho inclinado da
curva de pertinência (ou cruzar um limiar). Como a agregação é uma média geométrica ponderada, essa
queda se propaga de forma multiplicativa. O deslocamento climático é espacialmente suave
(Seção~\ref{sec:atlas-present}), mas o efeito em $\Delta S$ \emph{não} é uniforme. Um local perto de
uma fronteira de classe da FAO, por ter solo ou terreno marginais, desce de classe com um pequeno recuo
climático; um local com solo e terreno excelentes absorve o mesmo recuo sem mudar de classe. Daí a
realocação da cultura mais apta, concentrada em fronteiras específicas (Figura~\ref{fig:best-crop} do
Anexo), e não um deslocamento uniforme do mapa inteiro. Em suma: terreno e solo determinam
\emph{onde} o sistema é sensível; o clima fornece a variação \emph{temporal} que decide se essa
sensibilidade se materializa.

As mudanças médias territoriais vão de alguns décimos de milésimo, nas culturas comerciais
resilientes, a cerca de nove centésimos, na geração fotovoltaica sob a trajetória mais severa
(Tabela~\ref{tab:shift-mean}; Figura~\ref{fig:shift}). A projeção identifica três segmentos expostos,
nesta ordem: a geração fotovoltaica, a conservação e as outras culturas anuais. A ordenação é a mesma
em duas das três medidas disponíveis: o $\Delta S$ médio ($-0{,}091$, $-0{,}058$ e $-0{,}052$) e a
parcela de células rebaixadas de classe da FAO ($32{,}9\%$, $25{,}4\%$ e $21{,}8\%$). Na terceira, a
perda de área apta em pontos percentuais, a geração fotovoltaica continua destacada ($-15{,}8$), mas
as outras culturas anuais ($-13{,}0$) passam à frente da conservação ($-11{,}1$). Medida em termos
relativos à área apta de partida, a perda volta à ordem original ($-58{,}7\%$, $-28{,}9\%$ e
$-28{,}7\%$). A identificação do segmento mais exposto é, portanto, robusta ao indicador escolhido; a
ordem entre o segundo e o terceiro, não. Nos quatro segmentos restantes, o sinal é pequeno: soja,
cana-de-açúcar e pecuária ficam no piso de arredondamento ($|\Delta S| \le 0{,}0014$ na pior
combinação de cenário e janela), e a piscicultura melhora. A superfície de conservação inclui dois
fatores climáticos sensíveis (aridez e calor do trimestre mais quente); por isso, sua aptidão cai sob
aquecimento e ressecamento mais acentuadamente que a dos segmentos de produção.

\begin{table}[htbp]
\centering
\caption{Mudança média territorial projetada na aptidão, $\Delta S$ (futuro menos
presente), por segmento. Valores negativos indicam perda de aptidão.}
\label{tab:shift-mean}
\small
% GENERATED by tools/gen_thesis_tables.py -- do not edit by hand.
% Edit the source CSV and regenerate; see CLAUDE.md §9.
\begin{tabular}{@{}lrrrr@{}}
\toprule
 & \multicolumn{2}{c}{SSP2-4.5} & \multicolumn{2}{c}{SSP5-8.5} \\
\cmidrule(lr){2-3} \cmidrule(lr){4-5}
Segmento & 2031--50 & 2051--70 & 2031--50 & 2051--70 \\
\midrule
Soja & $-$0,0000 & $-$0,0001 & $-$0,0000 & $-$0,0014 \\
Cana-de-açúcar & $-$0,0002 & $-$0,0002 & $-$0,0009 & $-$0,0007 \\
Outras culturas anuais & $-$0,0117 & $-$0,0236 & $-$0,0174 & $-$0,0519 \\
Piscicultura & $+$0,0117 & $+$0,0123 & $+$0,0122 & $+$0,0124 \\
Pecuária & $-$0,0000 & $-$0,0000 & $-$0,0000 & $-$0,0010 \\
Conservação & $-$0,0326 & $-$0,0377 & $-$0,0391 & $-$0,0581 \\
Geração fotovoltaica & $-$0,0300 & $-$0,0462 & $-$0,0389 & $-$0,0905 \\
\bottomrule
\end{tabular}

\end{table}

Como as pertinências se baseiam em limiares, a medida agronomicamente legível é a mudança
na área apta (S2+, $S \ge 0{,}50$), apresentada na Tabela~\ref{tab:shift-share} como parcela da terra
avaliada de cada segmento. As parcelas do presente e do futuro são calculadas sobre a mesma linha de
base rederivada (presente $=$ futuro $- \Delta S$). Esse presente rederivado difere do atlas do
presente em menos de um décimo de ponto percentual, o que confirma que a linha de base se cancela no
cálculo de $\Delta S$.

\begin{table}[htbp]
\centering
\caption{Parcela da terra avaliada de cada segmento classificada como apta
(S2+, $S \ge 0{,}50$), presente $\rightarrow$ futuro.}
\label{tab:shift-share}
\small
% GENERATED by tools/gen_thesis_tables.py -- do not edit by hand.
% Edit the source CSV and regenerate; see CLAUDE.md §9.
\begin{tabular}{@{}lrrrrr@{}}
\toprule
 &  & \multicolumn{2}{c}{SSP2-4.5} & \multicolumn{2}{c}{SSP5-8.5} \\
\cmidrule(lr){3-4} \cmidrule(lr){5-6}
Segmento & Presente & 2031--50 & 2051--70 & 2031--50 & 2051--70 \\
\midrule
Soja & 55,8\% & 55,8\% & 55,8\% & 55,8\% & 55,7\% \\
Cana-de-açúcar & 42,1\% & 42,1\% & 42,1\% & 42,0\% & 42,0\% \\
Outras culturas anuais & 45,1\% & 41,4\% & 38,2\% & 39,7\% & 32,1\% \\
Piscicultura & 22,0\% & 24,5\% & 24,6\% & 24,6\% & 24,6\% \\
Pecuária & 76,0\% & 76,0\% & 76,0\% & 76,0\% & 76,0\% \\
Conservação & 38,5\% & 32,1\% & 31,2\% & 30,9\% & 27,3\% \\
Geração fotovoltaica & 26,9\% & 19,8\% & 16,3\% & 17,8\% & 11,1\% \\
\bottomrule
\end{tabular}

\end{table}

A geração fotovoltaica é o segmento mais exposto. Seu $\Delta S$ (Tabela~\ref{tab:shift-mean}) se
agrava de $-0{,}030$ (SSP2-4.5, 2031--2050) para $-0{,}091$ (SSP5-8.5, 2051--2070), a maior magnitude
entre os segmentos, e sua área apta cai para menos da metade, de 26,9\% para 11,1\%. A razão é que a
redução de desempenho dos painéis com a temperatura é um efeito \emph{direto} do calor, que nenhum
manejo amortece: à medida que as planícies quentes cruzam o limiar de calor, a superfície perde área
apta com facilidade.

A conservação é o segundo segmento mais exposto, com $\Delta S$ de $-0{,}033$ a $-0{,}058$ e a
área apta caindo de 38,5\% para 27,3\%. A vegetação nativa é o uso da terra menos
protegido do clima: os fatores climáticos dos segmentos de produção ficam sobre um amplo
platô ótimo, porque o manejo os mantém perto de seu ótimo térmico. Já os fatores da conservação ficam
na rampa ativa da curva, e o sinal de aquecimento e ressecamento se transmite diretamente em perda de
terra de alto valor de conservação.

As outras culturas anuais (milho, algodão e sorgo) são o grupo de culturas mais exposto. Sua
mudança se agrava monotonicamente de $-0{,}0117$ para $-0{,}0519$, e a área apta se contrai de 45,1\%
para 32,1\% (13,0 pontos percentuais), a maior perda de área apta entre as culturas, ainda que abaixo
dos 15,8 pontos da geração fotovoltaica. Suas pertinências favorecem a regularidade da chuva e uma
janela térmica limitada, ambas reduzidas pela temperatura crescente e pelo aumento dos graus-dia de
crescimento. Esses resultados são coerentes com as tendências observadas de aquecimento e
ressecamento na fronteira agrícola da transição entre a Amazônia oriental e o Cerrado entre 1981 e
2020 \citep[p.~4]{marengo_increased_2022}, embora aquele estudo seja observacional, e não uma projeção.

Soja, cana-de-açúcar e pecuária, por sua vez, são resilientes: seu $\Delta S$ é essencialmente
nulo, exceto sob a combinação mais quente, em que se torna levemente negativo ($-0{,}0014$,
$-0{,}0007$ e $-0{,}0010$, respectivamente); suas áreas aptas variam menos de dois décimos de ponto
percentual em todo o período. A \textbf{piscicultura é o único segmento que ganha}: seu $\Delta S$ é
positivo nas quatro combinações de cenário e janela ($+0{,}012$), e sua área apta sobe de 22,0\% para
24,6\%. O mecanismo é o mesmo que penaliza os demais, com o sinal trocado: o aquecimento empurra a
temperatura do trimestre mais quente \emph{em direção} ao ótimo da piscicultura, e não para fora dele.

Analisada por zona (Tabela~\ref{tab:zone-shift} do Anexo), a mudança nas culturas acentua o gradiente
agroambiental existente, em vez de deslocá-lo. As perdas se concentram de forma sistemática na baixada
quente e arenosa (Zona~4), a zona mais quente, mais seca e mais pobre em carbono orgânico do
território. Sob o cenário e a janela mais severos, a aptidão das outras culturas anuais cai 0,093
nessa baixada; a da soja, 0,004; e a da cana-de-açúcar, 0,002.

No planalto de cultivo (Zona~1), as mesmas culturas ficam praticamente inalteradas. Na escarpa de
conservação (Zona~2), onde as culturas já são inaptas, quase não há o que perder ($-0{,}012$ no pior
caso; é a zona menos penalizada das quatro). A baixada drenada (Zona~3) fica em posição intermediária,
com $-0{,}053$ nas outras culturas anuais. Assim, o sinal climático penaliza mais as culturas onde elas
já são marginais e menos onde já são favorecidas: a projeção concentra a aptidão nas melhores terras,
em vez de redistribuí-la.

A decomposição da mudança em movimentos entre as classes discretas da FAO confirma que a projeção age
quase inteiramente por rebaixamento da terra, uma classe de cada vez, e majoritariamente nos três
segmentos expostos (Tabela~\ref{tab:downgrade-share} do Anexo). Para soja, cana-de-açúcar e pecuária,
mesmo sob a combinação mais severa, menos de 0,9\% das células mudam de classe; na piscicultura,
nenhuma. Nos três segmentos expostos, a parcela rebaixada cresce de forma monotônica com o cenário e a
janela: na geração fotovoltaica, de 11,2\% (SSP2-4.5, 2031--2050) a 32,9\% (SSP5-8.5, 2051--2070); na
conservação, de 14,4\% a 25,4\%; e nas outras culturas anuais, de 4,8\% a 21,8\%. A matriz de transição
completa para as outras culturas anuais (Tabela~\ref{tab:transition-othercrops} do Anexo) mostra que
esse rebaixamento se concentra no topo da escala e que \emph{nenhuma} célula melhora de classe. Sob
SSP5-8.5, em 2051--2070, a classe S1 recua de 12,4\% para 7,3\% do território, cedendo 5,1 pontos
percentuais à classe S2, enquanto 12,7 pontos de S2 descem a S3 e 3,7 pontos de S3 passam a inaptos.

A ordenação por cenário e janela se comporta como o raciocínio físico exige. Para todo segmento
sensível ao clima, a mudança é mais pronunciada sob a trajetória de maior emissão e na janela mais
tardia: SSP5-8.5 em 2051--2070 é o caso mais severo em seis dos sete segmentos, e as duas janelas de
curto prazo diferem pouco. A exceção é a cana-de-açúcar, cujo $\Delta S$ é marginalmente mais
negativo em SSP5-8.5 2031--2050 ($-0{,}0009$) do que em 2051--2070 ($-0{,}0007$). Trata-se de uma
diferença de dois décimos de milésimo, dentro do piso de arredondamento do segmento, e não de uma
inversão de mecanismo.

\begin{figure}[htbp]
\centering
\includegraphics[width=\textwidth]{fig_4_4.png}
\caption{Mudança projetada na aptidão, $\Delta S$, sob SSP5-8.5, 2051--2070, para os três
segmentos expostos (geração fotovoltaica, conservação e outras culturas anuais) e uma cultura
comercial resiliente (soja). A realocação da cultura mais apta está na
Figura~\ref{fig:best-crop} do Anexo.}
\label{fig:shift}
\end{figure}

\section{Uso atual da terra e validação (QP2 e validação da QP1)}\label{sec:realized}

A fase final confronta o potencial modelado com o uso atual da terra e com uma medida
independente de produtividade. Três conjuntos de dados do catálogo entram aqui: uma
classificação anual de uso da terra a 30~m agregada à grade; as unidades administrativas de
segundo nível (247 municípios, um dos quais é o Distrito Federal); e os produtos de
produtividade MODIS MOD17 (NPP anual e, para o validador intracultura, GPP da estação de
crescimento).

\subsection{Composição do uso atual da terra}\label{sec:realized-composition}

O uso atual da terra de GO~e~DF é dominado por dois usos que, juntos, ocupam quatro quintos do
território (Tabela~\ref{tab:realized-composition}): pastagem cultivada e em mosaico (45,4\%) e
vegetação nativa (35,4\%). O cultivo anual ocupa cerca de um sexto do território, majoritariamente com
soja (13,4\%); a cana-de-açúcar (2,3\%) e o grupo de algodão, arroz e mosaico de culturas (0,8\%) são
bem menores. Esta última linha merece uma ressalva: a coleção do MapBiomas usada aqui não distingue
milho nem sorgo, ambos diluídos na classe de mosaico, e ela não é, portanto, o equivalente exato do
segmento de outras culturas anuais no uso atual (Seção~\ref{sec:limitations}).

A aquicultura medida pelo MapBiomas apresenta uma limitação grave: ela ocupa 0,0006\% do território,
cerca de 2~km\textsuperscript{2}, três ordens de grandeza abaixo da menor lavoura mapeada.
Não se trata de um zero por ausência de classe, e sim de uma área ocupada real, pequena demais para
ancorar limiares ou validar presença, ponto retomado como limitação na Seção~\ref{sec:limitations}.

\begin{table}[htbp]
\centering
\caption{Composição do uso atual da terra de GO~e~DF, como parcela do território. As
classes não modeladas agregam usos que o MapBiomas distingue, mas que este estudo não avalia
(urbano, mineração, água, silvicultura, culturas perenes e aquicultura).}
\label{tab:realized-composition}
% GENERATED by tools/gen_thesis_tables.py -- do not edit by hand.
% Edit the source CSV and regenerate; see CLAUDE.md §9.
\begin{tabular}{@{}lr@{}}
\toprule
Uso atual & Parcela do território \\
\midrule
Soja & 13,4\% \\
Cana-de-açúcar & 2,3\% \\
Outras culturas anuais & 0,8\% \\
Aquicultura & 0,0\% \\
Pastagem & 45,4\% \\
Vegetação nativa & 35,4\% \\
Outras classes (não modeladas) & 2,8\% \\
\midrule
\textbf{Total} & \textbf{100,0\%} \\
\bottomrule
\end{tabular}

\end{table}

\subsection{Potencial \emph{versus} uso atual (QP2)}\label{sec:pot-vs-real}

O zoneamento obtido separa o uso atual da terra de forma coerente. A composição do uso atual dentro
de cada zona biofísica (Tabela~\ref{tab:zone-realized} do Anexo, comentada zona a zona na
Seção~\ref{sec:zoning-results}) mostra três padrões. A soja e a cana-de-açúcar se concentram no
planalto de cultivo (Zona~1). A vegetação nativa chega a quase 70\% da escarpa íngreme de conservação
(Zona~2), justamente a zona cujo melhor uso modelado, comparativo e absoluto, é a conservação. E a
baixada quente e arenosa (Zona~4), identificada como marginal para culturas, é dominada por pastagem.
Essa é a corroboração independente mais forte disponível para a QP1: uma partição não supervisionada
de variáveis biofísicas reproduz a geografia observada de cultivo, pastejo e vegetação remanescente, e
a zona de conservação é precisamente onde o Cerrado nativo de fato persiste.

O cruzamento das classes da FAO com o uso atual fornece, por segmento, a parcela de todo o território
de GO~e~DF que é apta para esse segmento, mas não está atualmente sob esse uso
(Tabela~\ref{tab:pot-real-gap}). A maior lacuna é a das outras culturas anuais, com 44,3\% do
território, seguidas da soja (41,5\%) e da cana-de-açúcar (39,5\%). Como a pastagem é de longe o maior
uso atual, a maior parte dessa área disponível é pastagem existente, e não vegetação nativa.

\begin{table}[htbp]
\centering
\caption{Lacuna entre potencial e uso atual: parcela de todo o território de GO~e~DF que
é apta (S2+) para um segmento, mas não está atualmente sob esse uso.}
\label{tab:pot-real-gap}
% GENERATED by tools/gen_thesis_tables.py -- do not edit by hand.
% Edit the source CSV and regenerate; see CLAUDE.md §9.
\begin{tabular}{@{}lrr@{}}
\toprule
Segmento & Área (km\textsuperscript{2}) & Parcela do território \\
\midrule
Outras culturas anuais & 153.190 & 44,3\% \\
Soja & 143.523 & 41,5\% \\
Cana-de-açúcar & 136.530 & 39,5\% \\
Piscicultura & 75.651 & 21,9\% \\
\bottomrule
\end{tabular}

\end{table}

\subsection{Validação da aptidão do presente (QP1)}\label{sec:validation-present}

As superfícies baseadas em conhecimento foram submetidas a três verificações independentes, e
todas as métricas abaixo são calculadas sobre a metade do território que não participou da
calibração (Seção~\ref{sec:holdout}).

Duas grades distintas operam aqui e não devem ser confundidas. A grade da partição tem
$6 \times 6$ células, que formam blocos de $\approx$130~km dimensionados como descreve a
Seção~\ref{sec:holdout}, e define o que é calibração e o que é validação. A grade da amostragem tem
$8 \times 8$ células e é apenas o estrato que espalha os pontos de avaliação por toda a extensão de
GO~e~DF, em vez de concentrá-los onde cada cultura é mais densa. A primeira separa treino de teste; a
segunda apenas evita que a amostra seja geograficamente enviesada dentro de cada metade.

Confrontadas com a presença atual de culturas, as superfícies ordenam bem a presença acima do fundo
para duas das três culturas avaliadas (Tabela~\ref{tab:presence-auc}). A cultura com melhor AUC é a
\textbf{cana-de-açúcar} (0,906), seguida da soja (0,864); pelo índice de Boyce, a ordem se inverte,
com a soja em 0,976 contra 0,967 da cana-de-açúcar. Nas duas métricas, ambas as culturas são bem
discriminadas. Esta e as demais verificações de presença usam a amostra estratificada por segmento e
por espaço descrita na Seção~\ref{sec:zoning}, o que já mitiga parte do otimismo causado pela
autocorrelação espacial em métricas de presença sobre amostras uniformes.

O caso da cana-de-açúcar merece uma ressalva de leitura. Sua área cultivada
permanece concentrada em poucos municípios (Seção~\ref{sec:realized-composition}) e desequilibrada
entre as duas metades da partição, na razão de 4,07 (Tabela~\ref{tab:split-balance} do Anexo).
Assim, a AUC de 0,906 se apoia em 298 pixels de presença na metade de validação. O número é alto,
mas a amostra que o sustenta é pequena e geograficamente concentrada.

\begin{table}[htbp]
\centering
\caption{Discriminação por presença: superfícies de aptidão confrontadas com a presença atual
de culturas (AUC e índice de Boyce, amostra estratificada por espaço e segmento).}
\label{tab:presence-auc}
\small
% GENERATED by tools/gen_thesis_tables.py -- do not edit by hand.
% Edit the source CSV and regenerate; see CLAUDE.md §9.
\begin{tabular}{@{}lrrrrrr@{}}
\toprule
 & \multicolumn{3}{c}{AUC} & \multicolumn{2}{c}{Índice de Boyce} & Presenças \\
\cmidrule(lr){2-4} \cmidrule(lr){5-6} \cmidrule(lr){7-7}
Segmento & calib. & valid. & $\Delta$ & calib. & valid. & (valid.) \\
\midrule
Soja & 0,819 & 0,864 & $+$0,044 & 0,927 & 0,976 & 2.665 \\
Cana-de-açúcar & 0,814 & 0,906 & $+$0,092 & 1,000 & 0,967 & 298 \\
Outras culturas anuais & 0,712 & 0,678 & $-$0,035 & 0,914 & 0,313 & 152 \\
\bottomrule
\end{tabular}

\end{table}

A discriminação por presença responde a uma pergunta binária: a superfície separa os pixels onde a
cultura está daqueles onde ela não está? Nada diz, porém, sobre se a terra classificada como mais
apta é também a que mais produz. As duas verificações seguintes tratam dessa segunda pergunta. Ambas
usam produtos MOD17 do MODIS (Seção~\ref{sec:phaseb}), mas em escalas e com papéis distintos. A
primeira correlaciona a aptidão com a NPP anual agregada por município e serve apenas de contexto
produtivo. A segunda, mais exigente, compara a aptidão com a GPP da estação de crescimento dentro
dos pixels de cada cultura, e é a que mais se aproxima de uma validação de produtividade.

No nível municipal, usa-se o $\rho$ de Spearman entre a aptidão média municipal e a NPP anual sobre
terra manejada (Tabela~\ref{tab:spearman-npp} do Anexo). Como a NPP descreve o contexto de produção,
e não o rendimento, ela é lida apenas como \emph{contexto} grosseiro de produtividade, conforme
explicitado na Seção~\ref{sec:phaseb}. A correlação é estatisticamente
significativa ($p < 0{,}01$) para os sete segmentos e positiva para seis deles. Em ordem decrescente:
outras culturas anuais ($+0{,}54$), pecuária ($+0{,}51$), conservação ($+0{,}42$), geração
fotovoltaica ($+0{,}38$), cana-de-açúcar ($+0{,}29$) e soja ($+0{,}22$). A piscicultura é a única
com correlação negativa ($-0{,}30$).

Vale notar que as outras culturas anuais lideram aqui uma classificação em que ficam em último lugar
na validação por presença. As duas leituras não se contradizem. A discriminação por presença pergunta
se a superfície ordena \emph{onde} a cultura está, e falha fora da amostra para um segmento de área
cultivada mínima. A correlação municipal pergunta se a superfície ordena \emph{quanto} o território
produz, agregada aos 159 municípios com terra manejada na metade de validação, e aí o sinal é o mais
forte dos sete. É a única evidência quantitativa favorável de que a tese dispõe para esse segmento, e
é por isso que ela é citada explicitamente nas conclusões sobre as outras culturas anuais.

Essa mistura de sinais foi antecipada na Metodologia: a NPP anual mede a assimilação de carbono da
vegetação, e não o rendimento agronômico, de modo que os segmentos com mais biomassa em pé pontuam
mais alto. Isso é coerente com o comportamento documentado dos produtos NPP/GPP do MODIS, que, sem
viés global, superestimam a produtividade em locais pouco produtivos e a subestimam nos muito
produtivos, inclusive em lavouras fertilizadas, cuja eficiência de uso da luz é subestimada
\citep[pp.~286, 288--289]{turner_evaluation_2006}. É coerente também com a literatura sobre a conversão
de GPP em rendimento. Em um modelo de eficiência de uso da luz a 30~m em Montana (EUA), a produção
estimada a partir da GPP se correlaciona quase perfeitamente com a produção declarada por condado
($r = 0{,}96$), mas o rendimento estimado se correlaciona bem mais fracamente com o rendimento medido
no talhão ($r \approx 0{,}44$) \citep[p.~16]{he_regional_2018}. Por isso, a NPP é lida aqui apenas
como contexto, e não como validador.

A segunda verificação, de produtividade intracultura, avalia a GPP da estação de crescimento (outubro a março)
apenas nos pixels efetivamente cultivados de cada cultura e a correlaciona com a aptidão dessa cultura
\emph{dentro} da própria cultura. Isso remove as confusões causadas pelos ciclos de fenologia e pousio
de coberturas distintas. Na metade de validação, o gradiente da soja é $+0,526$
($p < 0{,}001$, $n = 7.243$), e o das outras culturas anuais, $+0,575$, monotônico
($p < 0{,}001$, $n = 5.466$): uma aptidão modelada mais alta está associada a produtividade
\emph{mais alta} dentro da própria cultura, e não apenas \emph{não mais baixa}. A cana-de-açúcar, por
sua vez, mostra um gradiente negativo de $-0{,}234$ ($p < 0{,}001$, $n = 2.071$), magnitude que não
se pode descartar como ruído e que aponta na direção oposta à da aptidão modelada. Com amostras dessa
ordem, o valor de $p$ é pouco informativo: os três gradientes são significativos, e o que distingue
os casos é o sinal e a magnitude.

Essas verificações de produtividade são \emph{secundárias}, e não tentativas independentes de
validar rendimento. Em conjunto, sustentam uma conclusão favorável ao modelo, com uma exceção
declarada: a soja e as outras culturas anuais o corroboram dentro da própria cultura, mas a
cana-de-açúcar o \emph{contradiz} justamente no segmento de melhor discriminação por presença.
As duas medições respondem a perguntas diferentes (uma distingue cana-de-açúcar de não cana; a outra
modela a produtividade entre canaviais), e a resposta da primeira não necessariamente acompanha a resposta da segunda. A
divergência, porém, é real e não se resolve com os dados desta tese. O mais provável é que seja um
artefato do validador: a GPP do MODIS subestima a eficiência de uso da luz em lavouras intensamente
manejadas (parágrafo anterior), e a cana-de-açúcar é a cultura mais intensamente manejada das três.
Essa é, porém, uma explicação plausível e não testada aqui. As métricas de presença
(Tabela~\ref{tab:presence-auc}), cujo alvo é mais próximo do que se quer validar, continuam sendo os
validadores primários onde a amostra as sustenta.

O \emph{proxy} de NPP anual também separa fortemente as zonas. Uma ANOVA de um fator entre as
quatro zonas, na metade de validação, é altamente significativa ($F = 1.417{,}3$,
$p < 0{,}001$), e as médias de grupo se ordenam com a altitude (Tabela~\ref{tab:zone-anova}). O
planalto alto e argiloso (Zona~1) e a escarpa de conservação (Zona~2) são os mais produtivos (0,835 e
0,827), a baixada drenada (Zona~3) fica em posição intermediária (0,669), e a baixada quente e arenosa
(Zona~4) é marcadamente a menos produtiva (0,583). O zoneamento capta, assim, um gradiente real de
produtividade observada, e não apenas um gradiente presente nas variáveis de entrada.

\begin{table}[htbp]
\centering
\caption{NPP média do MOD17 (unidades escalonadas) por zona, na metade de validação da
partição espacial. ANOVA de um fator: $F = 1.417{,}3$, $p < 0{,}001$.}
\label{tab:zone-anova}
% GENERATED by tools/gen_thesis_tables.py -- do not edit by hand.
% Edit the source CSV and regenerate; see CLAUDE.md §9.
\begin{tabular}{@{}lrr@{}}
\toprule
Zona & NPP média (MOD17) & $n$ \\
\midrule
Z1 & 0,668 & 2.001 \\
Z2 & 0,820 & 6.344 \\
Z3 & 0,648 & 5.749 \\
Z4 & 0,756 & 1.301 \\
\bottomrule
\end{tabular}

\end{table}

\textbf{Verificação cruzada orientada por dados.} Uma floresta aleatória foi treinada com a presença
atual de soja, usando o mesmo conjunto padronizado de atributos e uma amostra estratificada por espaço
e por segmento. O treino foi feito na metade de calibração e a avaliação, na metade de validação
($n = 8.376$); o resultado foi comparado célula a célula com a classificação baseada em
conhecimento. Nas faixas de \citet{landis_measurement_1977} apresentadas na
Seção~\ref{sec:sota-validacao}, a concordância é razoável, e não forte ($\kappa$ de Cohen
$= +0{,}292$). Na metade de validação, o mapa de conhecimento classifica 45,4\% dos pixels avaliados
como aptos (S2+), contra 18,3\% aos quais o classificador orientado por dados atribui alta
probabilidade de presença. As duas parcelas se referem a essa metade, e não ao domínio completo a que
se refere a parcela S2+ da Seção~\ref{sec:atlas-present}. Esta é a única das verificações desta seção
em que a diferença entre as duas metades é um contraste real de treino e teste, e não territorial: o
$\kappa$ sobe de 0,220 na metade de calibração para 0,292 na de validação.

O $\kappa$ de Cohen é definido como
\begin{equation*}
  \kappa = \frac{p_o - p_e}{1 - p_e},
\end{equation*}
em que $p_o$ é a proporção observada de concordância célula a célula entre as duas
classificações (conhecimento \emph{versus} floresta aleatória) e $p_e$ é a proporção de
concordância esperada apenas ao acaso, dada a prevalência marginal de cada classe em cada mapa.
Quando as duas classificações marcam áreas de tamanhos muito diferentes, como as duas
marginais acima, a própria diferença entre elas impõe um teto ao $\kappa$. Mesmo com a maior
concordância compatível com essas áreas, ele não passaria de $\kappa_M \approx 0{,}42$; o valor
observado corresponde, portanto, a cerca de 69\% do máximo permitido
\citep[pp.~42--43]{cohen_coefficient_1960}. Isso, por si só, já explica por que ele fica abaixo do que
a AUC sugeriria.

Esse $\kappa$ modesto não contradiz a AUC de 0,864 apresentada na Tabela~\ref{tab:presence-auc}:
são duas perguntas diferentes sobre o mesmo resultado. A AUC pergunta se os pixels em que a soja
é realmente plantada recebem, em geral, uma nota de aptidão mais alta do que pixels escolhidos ao
acaso. A resposta é sim, de forma consistente. Já o $\kappa$ pergunta algo mais rígido: as duas
superfícies concordam sobre \emph{qual} parte do território marcar como apta? Aqui a resposta é apenas
parcial, pela diferença de área discutida acima. Como a área declarada apta pelo conhecimento é duas
vezes e meia a do classificador, o $\kappa$ agregado deve ser lido com as concordâncias específicas
positiva e negativa em mente \citep[pp.~551, 557]{cicchetti_high_1990}, que esta análise não apresenta.
Ou seja, $\kappa = 0{,}292$ não quer dizer que o modelo erre \emph{onde} a soja tende a ocorrer
(isso já foi respondido pela AUC), e sim que a área que o conhecimento especializado chama de apta e a
área efetivamente ocupada pela soja ainda têm tamanhos bastante diferentes.

Essa divergência de área, mesmo moderada, é informativa. O nicho orientado por dados
permanece mais estreito que o envelope biofísico: boa parte da terra apta está sem uso.
Esse é, precisamente, o sinal de potencial \emph{versus} uso atual da QP2, já que terra
biofisicamente capaz é deixada ociosa por razões econômicas, de posse da terra, de logística ou de
inércia. 

Em suma: o mapa baseado em conhecimento responde ``onde a terra é capaz'', e a floresta aleatória
responde ``onde a soja é de fato plantada''. A divergência entre os dois é exatamente a lacuna
quantificada na Seção~\ref{sec:pot-vs-real}.

\subsection{Unidades municipais de decisão}\label{sec:municipal}

Agregando-se os resultados aos 247 municípios, obtém-se a visão relevante para a decisão
(Figura~\ref{fig:municipal}). Por área majoritária, os municípios se distribuem pelas quatro zonas de
forma desigual: o planalto de cultivo (Zona~1) é a zona mais comum, com 94 municípios; seguem a escarpa
de conservação (Zona~2, 70 municípios), a baixada quente e arenosa (Zona~4, 45 municípios) e a baixada
drenada (Zona~3, 38 municípios). A escarpa de conservação ocupa apenas 10,4\% do território, mas é a
zona majoritária de 70 municípios, quase um terço do total. É um relevo de bordas e encostas disperso
por todo o estado, de modo que a responsabilidade pela zona de maior valor de conservação está dividida
entre muitas prefeituras, nenhuma das quais concentra parcela significativa dela. A baixada drenada,
inversamente, é a zona majoritária do menor número de municípios, o que faz da vocação para
piscicultura a mais geograficamente concentrada das quatro.

A aptidão média municipal da soja tem mediana de 0,512 (faixa de 0,08 a 0,88). Vinte e seis
municípios alcançam uma média de nível S1 ($\ge 0{,}75$), e 134 dos 247 são, em média, aptos (S2+).
Sobrepondo-se a projeção de meados do século, identificam-se os municípios vulneráveis:
\textbf{137 dos 247} enfrentam um declínio projetado das outras culturas anuais mais acentuado que
$-0{,}05$ (média de $-0{,}082$ e, no pior caso, $-0{,}202$). É na baixada quente e arenosa (Zona~4) que
eles são proporcionalmente mais frequentes (39 dos seus 45 municípios estão sinalizados), e os cinco
mais expostos estão todos nela. Esse achado responde à QP3: o nicho de milho, algodão e sorgo se reduz
mais rápido no terreno mais baixo e já marginal.

Duas classificações (\emph{rankings}) municipais resumem a leitura relevante para a decisão
(Tabela~\ref{tab:municipal-rank}). Um escore de \textbf{oportunidade} (aptidão média da
soja $\times$ parcela apta, mas sem uso) destaca municípios com alto potencial e ampla
área disponível em terra já convertida (43 pontuam acima de 0,4). Os líderes combinam aptidão da
soja de 0,71 a 0,86 com parcelas sem uso de 0,65 a 0,81. Um \emph{ranking} de
\textbf{vulnerabilidade}, pelo declínio projetado das outras culturas anuais, é dominado por municípios
da baixada quente e remota e da planície arenosa do vale do Araguaia, a noroeste, liderados por
São Miguel do Araguaia.

\begin{table}[htbp]
\centering
\caption{\emph{Rankings} municipais: cinco primeiros por oportunidade (aptidão da soja
$\times$ parcela apta, mas sem uso) e por vulnerabilidade (mudança projetada $\Delta S$ das
outras culturas anuais, SSP5-8.5, 2051--2070).}
\label{tab:municipal-rank}
% GENERATED by tools/gen_thesis_tables.py -- do not edit by hand.
% Edit the source CSV and regenerate; see CLAUDE.md §9.
\begin{tabular}{@{}lr@{\hspace{1.5cm}}lr@{}}
\toprule
Oportunidade (município) & Escore & Vulnerabilidade (município) & $\Delta S$ \\
\midrule
Adelândia & 0,609 & São Miguel do Araguaia & $-$0,202 \\
Firminópolis & 0,584 & Mundo Novo & $-$0,164 \\
Sanclerlândia & 0,571 & Novo Planalto & $-$0,150 \\
São Luís de Montes Belos & 0,561 & Nova Crixás & $-$0,148 \\
Gouvelândia & 0,556 & Bonópolis & $-$0,132 \\
\bottomrule
\end{tabular}

\end{table}

A soja e as outras culturas anuais foram escolhidas para este par de mapas municipais por papéis
distintos e complementares, e não por uma preferência arbitrária. A soja é a cultura de maior aptidão
absoluta e maior estabilidade climática (Seção~\ref{sec:shift}); seu mapa de oportunidade lê a lacuna
entre potencial e uso atual onde ela é mais informativa. As outras culturas anuais, por sua vez, são o
grupo de culturas com a queda de aptidão projetada mais acentuada (Seção~\ref{sec:shift}), o que as
torna o sinal natural de vulnerabilidade municipal. O equivalente para os demais cinco segmentos está na
Figura~\ref{fig:all-segment-suit} do Anexo.

\begin{figure}[htbp]
\centering
\includegraphics[width=0.85\textwidth]{fig_4_9.png}
\caption{Aptidão média municipal da soja (acima) e municípios sinalizados como vulneráveis
pelo declínio projetado das outras culturas anuais (abaixo). O mapa de subutilização da soja
(potencial \emph{versus} uso atual) está na Figura~\ref{fig:underused-soy} do Anexo.}
\label{fig:municipal}
\end{figure}
